\section{三条 Attention 路线：Kimi、DeepSeek、MiniMax}

前面讲 Kimi，本章把三条路线放在一起比较。Kimi 走 Linear/Hybrid，DeepSeek 走 Sparse Attention，MiniMax M2 从 M1 的 Linear/Hybrid 回到 Full Attention。这三个选择都在回答同一个问题：怎样在长上下文 decoding 中降低成本，同时不牺牲太多表达能力和泛化能力。

\begin{figure}[H]
\centering
\includegraphics[width=0.96\textwidth]{linear-vs-sparse-vs-full.png}
\caption{三条 Attention 路线：Kimi、DeepSeek、MiniMax 的不同押注。自制概念图，依据 00:20:20--00:27:00 对谈内容整理。}
\end{figure}

\begin{knowledgebox}{读图：三条路线解决同一个痛点}
Linear/Hybrid 节省 KV cache，Sparse 减少每步关注 token，Full Attention 表达力强但贵。Kimi 选择混合线性，DeepSeek 更偏动态稀疏，MiniMax M2 回到全局注意力，说明不同团队对成本、稳定性和任务表现的权衡不同。
\end{knowledgebox}

\subsection{Kimi vs DeepSeek}

Kimi 的 hybrid 路线保留部分全局注意力层，大部分层换成线性注意力，因此能显著减少 KV cache，同时提高 decoding 效率。DeepSeek Sparse Attention 则通过 indexer 选择 Top-K token，减少每步 token 参与计算，但不一定减少每层 KV cache 存储。二者都面向长上下文效率，但优化点不同。

\begin{warningbox}{不要把 Linear 和 Sparse 当作谁绝对更先进}
Linear/Hybrid、Sparse 和 Full Attention 各有场景。表达能力、长程检索、KV cache、batch size、硬件实现、训练稳定性都会影响最终选择。模型公司回退到 Full Attention 不一定说明 Linear 失败，也可能说明当前任务和稳定性更需要全局表达。
\end{warningbox}

\subsection{MiniMax M2 为什么回到 Full Attention}

节目中提到，MiniMax M1 是大规模混合线性注意力的先驱之一，但 M2 又回到 Full Attention。一个可能原因是 Linear Attention 在一些 multi-hop reasoning 或复杂检索任务上表现不够稳；另一个原因是 Full Attention 的工程确定性、表达下限和训练经验仍然更成熟。架构选择不是只看推理速度，而是看综合表现。

\begin{knowledgebox}{三路线对照表}
\begin{tabularx}{\textwidth}{p{0.22\textwidth}p{0.25\textwidth}p{0.25\textwidth}X}
\toprule
路线 & 优点 & 代价 & 适合观察的问题 \\
\midrule
Full Attention & 表达力强、经验成熟 & KV cache 和长上下文成本高 & 多跳推理、复杂检索是否更稳。 \\
Linear/Hybrid & 省 KV cache、decoding 友好 & 需要弥补表达能力和长程交互 & 长 CoT、Agent 推理吞吐。 \\
Sparse Attention & 每步少看 token、硬件可协同 & 选 token 的 indexer 很关键 & 能否准确选中有用上下文。 \\
\bottomrule
\end{tabularx}
\end{knowledgebox}

\begin{importantbox}{架构选择是多目标优化}
模型公司不是在选“最优算法”，而是在性能、成本、稳定性、硬件实现、团队经验和产品场景之间做多目标优化。不同公司回到不同路线，是正常现象。
\end{importantbox}

\subsection{本章小结}

三条路线的分歧说明 Attention 还没有收敛。Kimi 更愿意用 hybrid 省成本，DeepSeek 更追求硬件协同下的动态稀疏，MiniMax 选择回到更稳的 Full Attention。未来可能出现组合方案。

